% Figure: the unit-system interface as a barrier vs a step. Two
% diagrams: one where the conversion is an audited interface, one
% where it is a memory operation inside a procedure.
\begin{figure}[htbp]
\centering
\begin{tikzpicture}[
  proc/.style={draw=engblack, rounded corners=2pt, minimum width=2.2cm,
    minimum height=0.9cm, align=center, font=\small},
  barrier/.style={draw=engred, thick, rounded corners=2pt,
    fill=engred!10, minimum width=2.5cm, minimum height=0.9cm,
    align=center, font=\small\bfseries},
  procstep/.style={draw=engblack!50, dashed, rounded corners=2pt,
    minimum width=2.2cm, minimum height=0.9cm, align=center,
    font=\small\itshape},
  arrow/.style={-{Stealth}, thick, draw=engblack!80},
  label/.style={font=\footnotesize\itshape, anchor=south},
]

% Top diagram: conversion as a barrier (good)
\node[proc, fill=engblue!10] (sysA1)  at (0, 3) {System A\\(imperial)};
\node[barrier]               (conv1)  at (3.5, 3) {audited\\conversion};
\node[proc, fill=englime!20] (sysB1)  at (7, 3) {System B\\(SI)};
\draw[arrow] (sysA1) -- (conv1);
\draw[arrow] (conv1) -- (sysB1);
\node[label] at (3.5, 4.0) {one location, documented, second pair of eyes};

% Bottom diagram: conversion as a step (bad)
\node[proc, fill=engblue!10] (sysA2)  at (0, 0.5) {System A\\(imperial)};
\node[procstep]                  (calc1)  at (2.3, 0.5) {value};
\node[procstep]                  (calc2)  at (4.6, 0.5) {operator\\mental conv.};
\node[procstep]                  (calc3)  at (7, 0.5)   {value};
\node[proc, fill=englime!20] (sysB2)  at (9.5, 0.5) {System B\\(SI)};
\draw[arrow] (sysA2) -- (calc1);
\draw[arrow] (calc1) -- (calc2);
\draw[arrow] (calc2) -- (calc3);
\draw[arrow] (calc3) -- (sysB2);
\node[label] at (4.6, -0.5) {multiple steps, undocumented, no cross-check};

\end{tikzpicture}
\caption[Unit-system interface: barrier vs step]{Two ways a unit
conversion can sit between two systems. Top: the conversion is an
explicit, audited interface (a barrier); the conversion happens in
one place that the engineering team owns and reviews. Bottom: the
conversion is a step inside the processing chain (memory work for
the operator); the chain has many opportunities for error and no
single place where the conversion can be checked. The Mars Climate
Orbiter and the Air Canada Flight 143 failures both involved
conversions as steps rather than barriers.}
\label{fig:vol01:ch01:unit-interface}
\end{figure}
